\documentclass[10pt]{article}

% ---------- Document Identity (update these only) ----------
\newcommand{\DocStage}{}
\newcommand{\DocTitle}{Your Road to Tier One SOF}
\newcommand{\ProgramName}{182-Week Civilian-to-Selection Readiness Pipeline}
\newcommand{\ProgramSubtitle}{}

% ---------- Page ----------
\usepackage[legalpaper,left=0.5in,right=0.5in,top=0.6in,bottom=0.45in]{geometry}

% ---------- Typography ----------
\usepackage[T1]{fontenc}
\usepackage[utf8]{inputenc}
\usepackage{libertinus}
\usepackage{microtype}
\usepackage{parskip}
\usepackage{ragged2e}
\usepackage{textcomp}
\AtBeginDocument{\color{black}}
\AtBeginDocument{\pagecolor{white}}
\AtBeginDocument{\justifying}

% ---------- Landscape pages ----------
\usepackage{pdflscape}

% ---------- Color ----------
\usepackage[table]{xcolor}
\definecolor{StageBlue}{HTML}{1F4E79}
\definecolor{StageBlueDark}{HTML}{163A5A}
\definecolor{LightGray}{HTML}{F2F4F7}
\definecolor{MidGray}{HTML}{D9DEE5}
\definecolor{RuleGray}{HTML}{C7CED8}
\definecolor{HeaderInk}{HTML}{000000}
\definecolor{Ink}{HTML}{000000}

% ---------- Utility ----------
\usepackage{enumitem}
\sloppy
\setlength{\emergencystretch}{2em}

% ---------- Boxes / UI ----------
\usepackage[most]{tcolorbox}

\tcbset{
  charterTitle/.style={
    enhanced,
    colback=StageBlue!4,
    colframe=StageBlue,
    boxrule=1.25pt,
    arc=3pt,
    left=16pt,right=16pt,top=14pt,bottom=14pt,
    borderline north={3.2pt}{0pt}{StageBlue},
    borderline west={4pt}{0pt}{StageBlue},
    borderline south={1.25pt}{0pt}{StageBlue},
    drop shadow={black!25!white},
  },
  charterRule/.style={
    enhanced,
    colback=LightGray,
    colframe=RuleGray,
    boxrule=0.85pt,
    arc=3pt,
    left=12pt,right=12pt,top=10pt,bottom=10pt,
    borderline west={3pt}{0pt}{StageBlue},
  },
  charterBadge/.style={
    enhanced,
    colback=StageBlue,
    coltext=white,
    colframe=StageBlueDark,
    boxrule=0.6pt,
    arc=2pt,
    left=6pt,right=6pt,top=3pt,bottom=3pt,
  }
}
\newtcbox{\Badge}{nobeforeafter,charterBadge}

% ---------- Lists ----------
\setlist[itemize]{leftmargin=1.5em, itemsep=0.25em, topsep=0.25em}
\setlist[enumerate]{leftmargin=1.8em, itemsep=0.25em, topsep=0.25em}

% ---------- TOC ----------
\usepackage{tocloft}
\setcounter{tocdepth}{3}
\setlength{\cftbeforesecskip}{3pt}
\setlength{\cftbeforesubsecskip}{2pt}
\setlength{\cftbeforesubsubsecskip}{0pt}
\renewcommand{\cftsecfont}{\bfseries}
\renewcommand{\cftsecpagefont}{\bfseries}
\renewcommand{\cftsubsecfont}{\normalfont}
\renewcommand{\cftsubsecpagefont}{\normalfont}
\renewcommand{\cftsubsubsecfont}{\normalfont}
\renewcommand{\cftsubsubsecpagefont}{\normalfont}
\setlength{\cftsecindent}{0pt}
\setlength{\cftsubsecindent}{1.25em}
\setlength{\cftsubsubsecindent}{2.8em}
\setlength{\cftsecnumwidth}{2.1em}
\setlength{\cftsubsecnumwidth}{2.8em}
\setlength{\cftsubsubsecnumwidth}{3.6em}

% Remove the built-in TOC heading so only the custom heading is shown
\makeatletter
\renewcommand{\tableofcontents}{\@starttoc{toc}}
\makeatother

% ---------- Header/Footer: page number only ----------
\usepackage{fancyhdr}
\pagestyle{fancy}
\fancyhf{}
\renewcommand{\headrulewidth}{0pt}
\renewcommand{\footrulewidth}{0pt}
\fancyfoot[C]{\thepage}

% ---------- Section formatting ----------
\usepackage{titlesec}

\titleformat{\section}
  {\Large\bfseries\color{Ink}}
  {\thesection}{0.6em}{}
  [\vspace{0.25em}\textcolor{StageBlue}{\rule{\linewidth}{0.8pt}}]

\titleformat{\subsection}
  {\large\bfseries\color{Ink}}
  {\thesubsection}{0.6em}{}

\titleformat{\subsubsection}
  {\normalsize\bfseries\color{Ink}}
  {\thesubsubsection}{0.6em}{}

\titlespacing*{\section}{0pt}{0.95em}{0.35em}
\titlespacing*{\subsection}{0pt}{0.65em}{0.25em}
\titlespacing*{\subsubsection}{0pt}{0.55em}{0.2em}

% ---------- Table engine ----------
\usepackage{tabularray}
\UseTblrLibrary{booktabs}

% ---------- tabularray: GLOBAL table treatment ----------
\SetTblrInner{
  cells = {fg=black, bg=white, valign=t},
}

% ---------- Header helpers ----------
\newcommand{\Hdr}[1]{\shortstack[c]{\strut #1\strut}}

% ---------- Lines ----------
\newcommand{\TitleLine}{\textcolor{StageBlue}{\rule{0.98\linewidth}{0.95pt}}}
\newcommand{\SubTitleLine}{\textcolor{RuleGray}{\rule{0.98\linewidth}{0.65pt}}}

% ---------- Continued on next page ----------
\newcommand{\ContinuedOnNextPageInline}{%
  \par
  \vfill
  \noindent\textcolor{RuleGray}{\rule{\linewidth}{1pt}}\vspace{-2pt}
  \noindent\hfill\textit{Continued on next page}\par\vspace{-10pt}
  \noindent\textcolor{RuleGray}{\rule{\linewidth}{1pt}}
  \clearpage
}

% ---------- PDF hyperlinks ----------
\usepackage[hidelinks]{hyperref}
\hypersetup{
  colorlinks=false,
  pdfborder={0 0 0},
  linktoc=all,
  pdftitle={\DocTitle},
  pdfauthor={\ProgramName}
}

% =========================================================
\begin{document}

\section*{7.12 — Recovery, Mobility, and Tissue-Care Implements}
\phantomsection
\addcontentsline{toc}{section}{7.12 — Recovery, Mobility, and Tissue-Care Implements}

Overview \textbf{ID}: 7.12\\
Overview Name: Recovery, Mobility, and Tissue-Care Implements\\
Version: v1.0\\
Governing Status: Draft — Non-Deployable\\
Scope Boundary Statement: Covers non-medical recovery tools, mobility-support implements, tissue-care implements, and comfort-support tools used to reduce execution friction; does not provide programming, diagnosis, treatment, or rehabilitation prescriptions.

\subsection*{7.12.1 Purpose \& Scope}
\phantomsection
\addcontentsline{toc}{subsection}{7.12.1 Purpose \& Scope}

This category covers:

\begin{itemize}
\item Mobility-support tools used to maintain or restore practical range of motion needed for safe movement mechanics across phases.
\item Tissue-care implements for self-massage and gentle release used to reduce soreness friction and support adherence.
\item Recovery comfort tools that reduce daily execution friction and help preserve consistency (especially in higher volume or higher consequence phases).
\end{itemize}

This overview crosswalks Stage 6 timing to the recovery and mobility support conditions required for safe execution and stable evidence posture.

This category does not cover:

\begin{itemize}
\item Sleep aids and ingestible sleep / stress aids: Stage 7.06.
\item Medical treatment, rehabilitation prescriptions, diagnosis, or “fixing injuries.”
\item Hygiene / sanitation supplies as a category: Stage 7.20 (cross-reference only).
\item Protective equipment: Stage 7.19 (cross-reference only when relevant).
\end{itemize}

This overview does not provide programming, test protocols, calendars, medical diagnosis, treatment, or rehabilitation prescriptions.

\subsection*{7.12.2 Governance And Safety Boundaries}
\phantomsection
\addcontentsline{toc}{subsection}{7.12.2 Governance And Safety Boundaries}

\begin{itemize}
\item Stage 0 boundary alignment: stop rules, safety override doctrine, conflict-resolution hierarchy, and evidence-admissibility controls apply to all 7.12 selections and substitutions.
\item Support-layer rule: these tools support consistency and comfort; they do not override stop rules or justify ignoring pain signals.
\item Non-medical posture: no diagnosis, no cure claims, and no rehab prescriptions.
\item Safe-use boundaries:
  \begin{itemize}
  \item avoid aggressive tissue work on acute pain/swelling,
  \item avoid targeting joints directly,
  \item discontinue and regress if symptoms worsen,
  \item stop immediately if sharp pain, numbness, or tingling occurs.
  \end{itemize}
\item Pressure/intensity safety rule (operational, mandatory):
  \begin{itemize}
  \item pressure must stay <=5/10 discomfort,
  \item stop on sharp pain or numbness/tingling.
  \end{itemize}
\item Change-control posture:
  \begin{itemize}
  \item introduce tools conservatively,
  \item do not add multiple novel recovery tools at once near gate windows to reduce confounding and to preserve comparability of recovery state.
  \end{itemize}
\item Tools are not substitutes for medical care:
  \begin{itemize}
  \item persistent or worsening pain, especially gait-altering pain, triggers evaluation posture per governance.
  \end{itemize}
\end{itemize}

\subsection*{7.12.3 Tier Doctrine}
\phantomsection
\addcontentsline{toc}{subsection}{7.12.3 Tier Doctrine}

\begin{itemize}
\item Price band convention: \$ = minimal household-cost item; \$\$ = modest purpose-built item; \$\$\$ = expanded-capability or redundancy item.
\item N/A rule: use N/A only when a field is genuinely not applicable and omission does not obscure safety, maintenance, or phase-timing logic.
\item Substitution posture: substitutions are acceptable only when they preserve low-risk comfort support, stability, and conservative pressure boundaries; substitutes that increase slip risk, pressure, or symptom aggravation break intent.
\item Tier 0: household-only posture using floor space + towel/pillow substitutes and a household ball for light tissue work; minimal capability; no purchases required.
\item Tier 1: minimum viable kit supporting daily session structure and deload compliance (mat, basic roller/ball, strap, gel pack).
\item Tier 2: expanded kit improving coverage and reducing friction (slant board/wedge, additional density options, optional additional small tools).
\item Tier 3: overbuilt posture emphasizing redundancy and selective specialized tools (massage gun is Tier 3 only; conservative-use boundaries), plus expanded heat/cold options.
\end{itemize}

\subsection*{7.12.4 Selection Criteria \& Fit Checks}
\phantomsection
\addcontentsline{toc}{subsection}{7.12.4 Selection Criteria \& Fit Checks}

Fit and safety checklist (operational):

\begin{itemize}
\item Cleanability and durability: tools must tolerate sweat and frequent use without degrading quickly.
\item Density/firmness matching tolerance:
  \begin{itemize}
  \item start conservative; obese beginner posture may require softer entry to avoid bruising and aversion.
  \end{itemize}
\item Stability and slip resistance:
  \begin{itemize}
  \item wedges/boards must be stable, non-slip, and anti-tip; if stability cannot be guaranteed, they are not authorized.
  \end{itemize}
\item Noise constraints:
  \begin{itemize}
  \item Motorized devices (massage gun) must consider apartment/noise limits; do not create compliance friction.
  \end{itemize}
\item Fit checks (operational):
  \begin{itemize}
  \item tool use does not provoke numbness/tingling,
  \item tool use does not provoke sharp pain,
  \item tool use does not increase next-day pain beyond baseline trend.
  \end{itemize}
\end{itemize}

\subsection*{7.12.5 Setup, Safety, and Anchoring}
\phantomsection
\addcontentsline{toc}{subsection}{7.12.5 Setup, Safety, and Anchoring}

Setup doctrine (no programming):

\begin{itemize}
\item Safe-use boundaries (mandatory):
  \begin{itemize}
  \item pressure <= 5/10 discomfort,
  \item avoid targeting joints directly,
  \item stop on sharp pain or numbness/tingling,
  \item avoid acute swollen areas.
  \end{itemize}
\item Environment controls:
  \begin{itemize}
  \item non-slip flooring and clear space; remove trip hazards.
  \item stable mat placement; wedge/board placed on stable surface.
  \end{itemize}
\item Hygiene posture:
  \begin{itemize}
  \item wipe-down after use; prevent shared-surface contamination.
  \end{itemize}
\item Change-control:
  \begin{itemize}
  \item do not introduce multiple new tools simultaneously near decision windows; stabilize earlier to reduce confounding.
  \end{itemize}
\end{itemize}

\subsection*{7.12.6 Maintenance, Replacement, and Storage}
\phantomsection
\addcontentsline{toc}{subsection}{7.12.6 Maintenance, Replacement, and Storage}

\begin{itemize}
\item Replacement cues:
  \begin{itemize}
  \item foam cracking, loss of rebound, sticky/degraded surfaces, or slip-risk from worn coatings.
  \end{itemize}
\item Cleaning cadence:
  \begin{itemize}
  \item wipe-down after use; deeper clean weekly.
  \item store dry to prevent mold odor and material breakdown.
  \end{itemize}
\item Storage:
  \begin{itemize}
  \item dry, out of direct sunlight; avoid compressing foam tools permanently.
  \item keep floor space clear to reduce trip hazards.
  \end{itemize}
\end{itemize}

\subsection*{7.12.7 Substitution Rules — Household Alternatives}
\phantomsection
\addcontentsline{toc}{subsection}{7.12.7 Substitution Rules — Household Alternatives}

\begin{itemize}
\item If no foam roller: towel roll or firm pillow substitute; keep pressure <=5/10.
\item If no mobility strap: towel or belt substitute with safe handling; avoid aggressive pulling.
\item If no gel pack: controlled warm shower or towel-based compress; comfort-only.
\item Documentation posture:
  \begin{itemize}
  \item Tools are optional supports; missing tools do not invalidate a session by themselves.
  \item if a tool absence causes movement regression or increases pain, regress exposure and prioritize safety; persistent limitations trigger evaluation posture rather than tool escalation.
  \end{itemize}
\end{itemize}

\begin{landscape}

\subsection*{7.12.8 Phase Integration Map — Required}
\phantomsection
\addcontentsline{toc}{subsection}{7.12.8 Phase Integration Map — Required}

\begingroup
\scriptsize
\setlength{\tabcolsep}{2pt}
\renewcommand{\arraystretch}{1.12}

\begin{longtblr}{
  width=\linewidth,
  rowhead=1,
  colspec = {
    Q[l,wd=1.05in]
    Q[l,wd=1.55in]
    X[l,2.75]
    X[l,3.65]
  },
  hlines,
  vlines,
  cells = {hsep=6pt, vsep=6pt, halign=l, valign=t},
  row{1} = {bg=StageBlue!15, fg=black, font=\bfseries, halign=l, valign=m},
  row{odd} = {bg=white},
  row{even} = {bg=LightGray},
}
\Hdr{Phase} &
\Hdr{Required Tier (from Stage 6)} &
\Hdr{What Changes / Becomes Mandatory} &
\Hdr{Notes} \\
Phase 00 &
Tier 1 &
Tier 1 minimum recovery / mobility kit required by Phase 00 (End): mat + basic roller/ball + strap + gel pack &
Phase 00 (End) is the first required point per Stage 6. Tools support adherence and mobility maintenance; they do not override stop posture. \\
Phase 01 &
Tier 1 &
Maintain Tier 1; increase consistency and cleaning discipline as volume rises &
Phase 01 (End): avoid introducing new tools inside Deload+Test windows; stabilize earlier. \\
Phase 02 &
Tier 1 &
Maintain Tier 1; strength exposure increases makes mobility consistency more consequential &
Phase 02 (End): prioritize shoulder/hip/ankle mobility maintenance without aggressive tissue work. \\
Phase 03 &
Tier 1 &
Maintain Tier 1; ruck capability becomes Tier 1 and increases foot/ankle/hip mobility relevance &
Phase 03 (End): higher consequence for ankle/calf mobility; avoid adding new tools in protected windows. \\
Phase 04 &
Tier 1 &
Maintain Tier 1; load-bearing increases friction and soreness; gel pack comfort use remains conservative &
Phase 04 (End / Mid): keep tools stable; do not escalate pressure to “fix” soreness. \\
Phase 05 &
Tier 1 &
Maintain Tier 1; endurance volume increases soreness friction; stability matters &
Phase 05 (End): use tools to reduce adherence friction, not to increase training dose. \\
Phase 06 &
Tier 2 &
Tier 2 becomes mandatory by Phase 06 (End): wedge/slant board option and expanded density/coverage &
Phase 06 (End) is first Tier 2 point per Stage 6. Emphasis is stability and repeatability under higher consequence phases. \\
Phase 07 &
Tier 2 &
Maintain Tier 2; hybrid density increases recovery importance &
Phase 07 (End): keep tool set stable; prevent overuse through conservative use boundaries. \\
Phase 08 &
Tier 2 &
Maintain Tier 2; load tolerance increases mobility and tissue-care operational importance &
Phase 08 (End): ankle/hip mobility and foot tissue-care boundaries become more operationally critical; still no aggressive pressure. \\
Phase 09 &
Tier 2 &
Maintain Tier 2; high \textbf{CAL} volume increases shoulder/elbow tissue-care relevance &
Phase 09 (End): use conservative tools to reduce friction; do not “roll through” pain. \\
Phase 10 &
Tier 2 &
Maintain Tier 2; multi-domain gates increase consequence of recovery collapse &
Phase 10 (End): stable tools and stable routines reduce variance; no novelty near protected windows. \\
Phase 11 &
Tier 2 &
Maintain Tier 2; recovery sensitivity increases consequence &
Phase 11 (End): prioritize consistency; reslot rather than forcing high-consequence weeks under unstable recovery. \\
Phase 12 &
Tier 3 &
Tier 3 becomes mandatory by Phase 12 (End): redundancy posture; massage gun optional Tier 3 only &
Phase 12 (End): Stage 6 indicates Tier 3 for reliability. Massage gun remains Tier 3 only with conservative-use boundaries; do not treat as required. \\
Phase 13 &
Tier 3 &
Maintain Tier 3; consolidation requires low variance and reliable recovery supports &
Phase 13 (End): redundancy exists to prevent compromised windows due to missing tools; do not introduce novelty near final gates. Massage gun remains optional and conservative-use boundaries still apply. \\
\end{longtblr}

\endgroup

\end{landscape}

7.12.8 Notes:

\begin{itemize}
\item If any upstream artifact requires wedge/slant board earlier than Phase 06 (End) without Stage 6 alignment, Requires Stage 8 review.
\item If any upstream artifact implies Massage Gun is required, Requires Stage 8 review.
\end{itemize}

\subsection*{7.12.9 Risks, Contraindications, and Red Flags}
\phantomsection
\addcontentsline{toc}{subsection}{7.12.9 Risks, Contraindications, and Red Flags}

\begin{itemize}
\item Overpressure risk:
  \begin{itemize}
  \item bruising and symptom aggravation occurs when pressure exceeds tolerance; enforce <= 5/10 discomfort rule.
  \end{itemize}
\item Masking risk:
  \begin{itemize}
  \item aggressive rolling/massage can mask pain signals and encourage unsafe exposure; tools must not override stop posture.
  \end{itemize}
\item Red flags (stop tool use and follow escalation posture):
  \begin{itemize}
  \item swelling, sharp pain, numbness/tingling, worsening symptoms.
  \item persistent gait-altering pain → regress exposure and seek evaluation posture.
  \end{itemize}
\item Tool hygiene risks:
  \begin{itemize}
  \item dirty tools and damp storage increase irritation risk; clean routinely and store dry.
  \end{itemize}
\item Requires Stage 8 review triggers:
  \begin{itemize}
  \item any mismatch where a phase demands Tier 1 recovery tools earlier than Stage 6,
  \item any attempt to move Massage Gun into Tier 1–2 or to frame it as required,
  \item any attempt to introduce aggressive "injury-fix" claims inside Stage 7.12.
  \end{itemize}
\end{itemize}

\end{document}